\cotefile{162-5}
\scanpage{27}
5.\quad Les $\Theta$-Algèbres loc libres \inserted{de t.f. $\mathscr T$} sur $X_0$ correspondent exactement aux complexes de De-Rham sur $X_0$ construits sur des $\Omega$ loc. libres de t. f, par les correspondances de $3^\circ$. et $4^\circ$.

6.\quad Revenons au cas d'une catégorie formelle sur $X_0$, définie par $A=(A_\alpha)$ sur $X_0\deleted{\unsure{\mathcal C}}$\,...., et le complexe de De Rham associé $(\Omega^*\unsure{\underset{\nearrow}{=}|=}\bigwedge^*\Omega,\delta)$.

\uline{Th. de Quillen} \inserted{Supposons $X_0$ de car. nulle} Le foncteur $A\longmapsto(\Omega^*,\delta)$ induit une \uline{équivalence} entre la catégorie des catégories formelles \inserted{\unsure{adiques|-adiques}} sur $X_0$ dont le $\Omega\to$ \uline{plat}, et la catégorie des \deleted{Algèbres}\inserted{complexes} de De Rham $(\bigwedge^\alpha\Omega,\delta)$ sur $X_0$ dont le $\Omega$ \unsure{est|soit} \uuline{plat.} [Les $A$ en question ont la vertu que $\mathrm{Sym}\,\Omega\longrightarrow\mathrm{Gr}(A)$ est un isomorphisme, ce qui s'apparente \unsure{à|:} une propriété de ``lissité différentielle''.

Ceci contient le th. de Cartier sur les groupes formels en car. nulle.

Il y a une variante ``à puiss. divisées'', quand on met sur l'idéal d'augmentation de $A$ \inserted{$=(A_n)_{n\in\mathbb N}$} une structure d'\deleted{Algèbre} puiss divisées, qui passe au quotient sur les $A_n$, chaque $I_n$ étant nilpotent d'ordre $n$, et la filtration étant celle provenant des puiss. divisées. [À voir si dans cet énoncé il faut toujours des \unsure{relat.} entre la structure \deleted{\unsure{de}} puiss. divisées, \deleted{et} l'application $\Delta$ \ldots]

\margin{26}
